\chapter{Production Proof -- Wie AI Engineers bewertet werden}
\label{chap:bewerbungsprozess}

% ============================================================
% LERNZIELE
% ============================================================
\begin{lernziele}
Nach diesem Kapitel kannst du:
\begin{itemize}
    \item Ein AI-System gegen die sechs Produktionsdimensionen evaluieren (Latenz, Kosten, Zuverlässigkeit, Halluzination, Eval, Observability)
    \item Portfolio-Projekte als Production Proof strukturieren (Monitoring, Eval-Loop, Fallback, Caching, Deployment)
    \item System-Design-Interviews als Debugging-Probleme lösen (kaputte Architektur reparieren)
    \item Die eigenen Systeme gegen die Metriken messen, die in Produktion zählen
    \item Failure-Mode-Analysen für RAG-, Agent- und Chat-Systeme durchführen
\end{itemize}
\end{lernziele}

% ============================================================
% MOTIVATION
% ============================================================
\section{Motivation}

Dieses Buch hat ein komplettes AI-Engineering-Spektrum abgedeckt: von der Modell-Wahl über RAG-Architektur und Agent-Orchestrierung bis zur Inference-Optimierung und Sicherheit. Aber das Buch endet nicht mit dem Code -- es endet mit der Frage, die jedes Produktionsteam stellt: \glqq Funktioniert dieses System unter echtem Druck?\grqq

In der Praxis werden AI Engineers nicht danach bewertet, ob sie ein Konzept erklären oder ein Notebook schreiben können. Sie werden bewertet an den Systemen, die sie gebaut haben und die in Produktion laufen. Ein RAG-System mit 99,5\,\% Availability, 2-Sekunden-P95-Latenz und 0,002~EUR Cost per Query sagt mehr über eine Ingenieurin aus als jedes Zertifikat.

Dieses Kapitel reframt den gesamten Bewertungsprozess: weg vom Career-Coaching, hin zur System-Analyse. Interviews sind Produktions-Audits. Portfolios sind Stress-Test-Ergebnisse. Und die entscheidende Frage ist nicht \glqq Was weisst du?\grqq{} sondern \glqq Was hast du gebaut -- und hält das Stand?\grqq

% ============================================================
% SYSTEM EVALUATION
% ============================================================
\section{Hiring als System Review}

Ein technisches Interview für AI Engineering ist kein Quiz. Es ist ein \textbf{Production System Review}. Der Interviewer hat ein Bild eines funktionierenden Systems im Kopf und prüft, ob deine Architektur-Entscheidungen diesen Bild entsprechen. Die Bewertungskriterien sind identisch mit denen eines Post-Mortems oder einer Architecture-Review:

\subsection{Die sechs Produktionsdimensionen}

\begin{table}[H]
\centering
\begin{tabularx}{\linewidth}{lXl}
\toprule
\textbf{Dimension} & \textbf{Fragestellung} & \textbf{Metrik} \\
\midrule
Latenz & Wie schnell antwortet das System unter Last? & P95 TTFT $<$ 500ms \\
Kosten & Was kostet ein einzelner Request im Durchschnitt? & EUR/Query $<$ 0,005 \\
Zuverlässigkeit & Wie oft fällt das System aus? & Uptime $>$ 99,5\,\% \\
Halluzination & Wie oft erfindet das System Fakten? & Hall. Rate $<$ 2\,\% \\
Eval-Pipeline & Kannst du Regressionen automatisch erkennen? & CI-Pass-Rate $>$ 90\,\% \\
Observability & Kannst du einen Fehler in 5 Minuten lokalisieren? & MTTR $<$ 30 Min \\
\bottomrule
\end{tabularx}
\caption{Die sechs Produktionsdimensionen -- gemessen in jedem System-Review}
\label{tab:prod-dims}
\end{table}

Jede Frage in einem Interview testet implizit eine dieser Dimensionen:

\begin{verbatim}
   "Design a customer support bot"
   -> Latenz: Wie verhinderst du, dass 10.000 gleichzeitige User
      das System überlasten?
   -> Kosten: Wann setzt du GPT-4o-mini ein, wann teurere Modelle?
   -> Zuverlaessigkeit: Was passiert, wenn die Vector-DB ausfaellt?
   -> Halluzination: Wie stellst du sicher, dass der Bot nur
      korrekte Produktinformationen liefert?
   -> Eval: Wie testest du 200 Varianten von Support-Fragen?
   -> Observability: Wie findest du heraus, warum ein User
      eine falsche Antwort bekommen hat?
\end{verbatim}

\subsection{Die System-Review-Checkliste}

Jede Architektur, die du in einem Interview präsentierst oder in einem Portfolio zeigst, wird gegen diese Checkliste geprüft:

\begin{enumerate}[leftmargin=*]
    \item \textbf{Latenz-Budget}: Hast du eine TTFT-Obergrenze definiert? Welche Optimierung (Caching, Streaming, kleineres Modell) setzt du unterhalb dieser Grenze ein?
    \item \textbf{Cost-Governance}: Hast du ein Budget pro Query und pro User? Wo sitzt der Cost-Controller in der Architektur?
    \item \textbf{Fallback-Kette}: Was passiert bei jedem einzelnen Ausfall -- Provider-Down, DB-Timeout, Embedding-Fehler, Rate-Limit?
    \item \textbf{Eval-Schleife}: Wie stellst du sicher, dass ein Prompt-Update das System nicht verschlechtert? Wo läuft das Golden Dataset?
    \item \textbf{Observability}: Welche Metriken loggst du? Wie sieht dein Dashboard aus? Wie findest du eine Regression in 5 Minuten?
    \item \textbf{Sicherheit}: Wo sitzt der Input-Guard? Wie verhinderst du Prompt Injection durch Tool-Ergebnisse?
\end{enumerate}

\section{Das Portfolio als Production Proof}

Ein Portfolio-Projekt wird nicht nach Code-Qualität allein bewertet. Es wird danach bewertet, ob es \textit{in Produktion überleben würde}. Ein Repository ohne Monitoring, Eval-Loop, Fallback-Strategien und Deployment-Konfiguration ist kein Production Proof -- es ist eine Übungsaufgabe.

\subsection{Die fünf Layer eines Production-Proof-Projekts}

\begin{verbatim}
   +--------------------------------------------------------------+
   |                Production-Proof-Projekt                       |
   |                                                              |
   |  Layer 5: Deployment (Docker, CI/CD, Cloud)                  |
   |  Layer 4: Observability (Tracing, Logging, Metriken)         |
   |  Layer 3: Eval-Loop (Golden Dataset, CI-Check)              |
   |  Layer 2: Fallback-Strategie (Provider, Cache, Degrade)      |
   |  Layer 1: Kostenkontrolle (Budget, Caching, Model-Routing)   |
   +--------------------------------------------------------------+
\end{verbatim}

\subsubsection{Layer 1 -- Kostenkontrolle}

Der auffälligste Unterschied zwischen einem Hobby-Projekt und einem Produktionssystem: Kostenbewusstsein.

\begin{lstlisting}[language=Python, caption=Cost-Awareness im Portfolio nachweisen]
class ProductionAwarePipeline:
    """Zeigt, dass Kosten, Latenz und Fehler behandelt werden."""

    def __init__(self):
        self.cache = SemanticCache(threshold=0.92)
        self.router = SemanticRouter()
        self.budget = BudgetController()
        self.metrics = MetricsCollector()

    async def handle(self, query: str, user_id: str) -> dict:
        start = time.monotonic()

        # 1. Budget-Check
        if not self.budget.check_request(user_id, 0.005):
            return {"error": "Budget limit reached",
                    "status": 429}

        # 2. Cache-Lookup (0-Latenz bei Hit)
        cached, score = self.cache.lookup(query)
        if cached:
            self.metrics.record(user_id, "cache_hit", 0.0)
            return {"response": cached, "cached": True}

        # 3. Intent-Routing
        route = self.router.route(query)
        response = await call_provider(route["model"], query)

        # 4. Kosten tracken
        cost = calculate_cost(response, route["model"])
        self.budget.record_cost(user_id, cost)
        self.metrics.record(user_id, "llm_call", cost)

        # 5. Cache-schreiben (je nach Intent)
        if route["intent"] not in ("personal", "sensitive"):
            self.cache.store(query, response)

        elapsed = time.monotonic() - start
        self.metrics.record_latency(user_id, elapsed)

        return {"response": response, "cached": False,
                "latency_ms": int(elapsed * 1000)}
\end{lstlisting}

\subsubsection{Layer 2 -- Fallback-Strategie}

Jedes Produktionssystem hat Ausfälle. Ein Production-Proof-Projekt zeigt, wie es damit umgeht:

\begin{lstlisting}[language=Python, caption=Fallback-Chain für LLM-Produktionssystem]
class FallbackChain:
    """Mehrstufige Fallback-Strategie fuer LLM-Calls."""

    def __init__(self):
        self.providers = [
            {"name": "primary", "model": "claude-sonnet-4",
             "client": anthropic.Client()},
            {"name": "backup", "model": "gpt-4o",
             "client": openai.Client()},
            {"name": "fallback", "model": "gpt-4o-mini",
             "client": openai.Client()},
        ]
        self.cache_ttl = 300

    async def execute(self, prompt: str,
                      timeout: float = 10.0) -> dict:
        errors = []

        for provider in self.providers:
            try:
                response = await asyncio.wait_for(
                    provider["client"].complete(prompt),
                    timeout=timeout,
                )
                return {
                    "response": response,
                    "provider": provider["name"],
                    "model": provider["model"],
                    "fallback_used": len(errors) > 0,
                }
            except Exception as e:
                errors.append({
                    "provider": provider["name"],
                    "error": str(e),
                })
                continue

        # Alle Provider ausgefallen -> Degraded Mode
        return {
            "response": "Service temporarily unavailable. "
                        "Please try again.",
            "provider": None,
            "model": None,
            "errors": errors,
            "degraded": True,
        }
\end{lstlisting}

\subsubsection{Layer 3 -- Eval-Loop}

Ein Projekt ohne Evaluation ist kein Produktionssystem. Zeige, dass du dein System systematisch testest:

\begin{verbatim}
   Eval-Loop im Portfolio:

   [Code Change] --> [Golden Dataset Run] --> [Score Check]
                         |                         |
                   [100+ Test Cases]         [Threshold: 85%]
                         |                         |
                    [LLM-Judge]          [Fail -> Block Merge]
                                              [Pass -> Deploy]
\end{verbatim}

\subsubsection{Layer 4 -- Observability}

Zeige Tracing, Metriken und Logging. Ein Interviewer will sehen, dass du einen Produktionsfehler in Minuten finden kannst:

\begin{lstlisting}[language=Python, caption=Observability-Layer im Portfolio]
class ObservabilityLayer:
    """Zeigt, dass Monitoring und Debugging implementiert sind."""

    def instrument(self, app):
        @app.middleware("http")
        async def trace_request(request, call_next):
            trace_id = str(uuid.uuid4())
            request.state.trace_id = trace_id

            # Startzeit
            start = time.monotonic()

            try:
                response = await call_next(request)
            except Exception as e:
                # Fehler mit Kontext loggen
                logger.error({
                    "trace_id": trace_id,
                    "error": str(e),
                    "path": request.url.path,
                    "user": request.headers.get("x-user-id"),
                })
                raise

            # Latenz tracken
            latency = time.monotonic() - start
            METRICS_HISTOGRAM.labels(
                path=request.url.path,
                status=response.status_code,
            ).observe(latency)

            # Response loggen (PII-free)
            logger.info({
                "trace_id": trace_id,
                "latency_ms": int(latency * 1000),
                "status": response.status_code,
                "tokens": response.headers.get("x-tokens-used"),
            })

            return response
\end{lstlisting}

\subsubsection{Layer 5 -- Deployment}

Jedes Projekt braucht \texttt{docker-compose.yml}, CI/CD und eine klare Startup-Anleitung. Eine einzelne \texttt{README.md}-Zeile \glqq pip install -r requirements.txt\grqq{} ist kein Deployment-Nachweis.

\begin{lstlisting}[language=bash, caption=Minimale Deployment-Konfiguration im Portfolio]
# docker-compose.yml -- Zeigt Infrastruktur-Denken
version: "3.8"
services:
  api:
    build: .
    ports: ["8000:8000"]
    env_file: .env
    depends_on:
      - redis
      - qdrant
    healthcheck:
      test: ["CMD", "curl", "-f", "http://localhost:8000/health"]
      interval: 30s

  redis:
    image: redis:7-alpine
    ports: ["6379:6379"]

  qdrant:
    image: qdrant/qdrant:latest
    ports: ["6333:6333"]
    volumes: ["./qdrant_storage:/storage"]
\end{lstlisting}

% ===========================================================%
% INTERVIEW ALS DEBUGGING
% ============================================================
\section{System Design als Debugging-Problem}

Ein System-Design-Interview testet nicht, ob du die perfekte Architektur aus dem Stegreif entwerfen kannst. Es testet, ob du die typischen Fehler erkennst, die in jedem AI-System auftreten. Der effektivste Ansatz: Behandle die Frage als kaputtes System, das du reparieren musst.

\subsection{Failure-Mode-Szenarien}

\subsubsection{Szenario A: Der Support-Bot halluziniert Produkte}

\textbf{Symptom:} Ein KI-Support-Bot für ein E-Commerce-System mit 500.000 Produkten hat eine Halluzinationsrate von 12\,\%. Kunden beschweren sich über erfundene Produkteigenschaften.

\textbf{Diagnose:}

\begin{verbatim}
   Moegliche Ursachen (geordnet nach Wahrscheinlichkeit):

   1. Retrieval-Qualitaet: 12% Halluzination = Embedding-Modell
      findet in 12% der Faelle irrelevante Chunks
      -> Check: Chunking-Strategie, Embedding-Modell, Top-K

   2. Prompt-Struktur: Der System-Prompt sagt "wenn du die Antwort
      nicht weisst, sage es dem User" -> wird ignoriert
      -> Check: Prompt-Testing, Few-Shot-Beispiele fuer "weiss nicht"

   3. Hybrid-Search: Nur Vector-Suche, keine Keyword-Suche
      -> Check: Bei exakten Produktnamen versagt die semantische Suche

   4. Kein Eval-Loop: Kein Golden Dataset -> keine Regression erkannt
\end{verbatim}

\textbf{Lösung:} Chunking auf 256 Tokens reduzieren, Cross-Encoder Re-Ranking einführen, Hybrid Search (Vector + BM25) implementieren, Golden Dataset mit 500 Fällen aufbauen, Eval-Check in CI/CD.

\subsubsection{Szenario B: Der Code-Review-Agent blockiert jeden PR}

\textbf{Symptom:} Ein KI-Code-Review-Agent markiert 40\,\% aller PRs als fehlerhaft. Das Team ignoriert die KI-Reviews komplett.

\textbf{Diagnose:}

\begin{verbatim}
   Drei unabhaengige Probleme:

   1. False-Positive-Rate zu hoch: Der Agent hat ein zu sensitives
      Modell (Halluziniert Sicherheitsprobleme)
      -> Loesung: Threshold erhoehen, zweite Meinung einholen

   2. Keine Learning-Loop: Falsche Markierungen werden nicht
      korrigiert -> das Modell lernt nicht aus Fehlern
      -> Loesung: Feedback-Button, Dataset erweitern

   3. Latenz: 4 Minuten Review-Zeit -> Entwickler warten nicht
      -> Loesung: Streaming-Response, parallele Analyse
\end{verbatim}

\subsubsection{Szenario C: Die Kosten eines RAG-Systems wachsen unkontrolliert}

\textbf{Symptom:} Ein RAG-System mit 10.000 täglichen Usern kostet plötzlich das Dreifache. Der CFO verlangt eine Erklärung.

\textbf{Diagnose:}

\begin{verbatim}
   Kosten-Treiber (nacheinander pruefen):

   1. Token-Verbrauch gestiegen: Neue Produktbeschreibungen
      sind 3x langer -> mehr Input-Tokens pro Query
      -> Check: Avg Input-Tokens/Query ueber Zeit

   2. Cache-Invalidierung: System-Prompt geandert -> Prompt-Cache
      invalid -> 100% Prefill-Kosten statt 10%
      -> Check: Cache-Hit-Rate ueber Zeit

   3. Modell-Wechsel (heimlich): Provider hat das Modell
      hinter "gpt-4o" auf eine teurere Version umgestellt
      -> Check: Model-String im Log

   4. Agent-Loop: Ein Multi-Agent-System ruft 8x statt 3x
      pro Task das LLM auf
      -> Check: Avg Steps/Task ueber Zeit
\end{verbatim}

\section{Technische Interview-Tiefe}

Die folgenden Themen werden in Senior-AI-Engineering-Interviews systematisch abgefragt. Nicht als Wissensfragen, sondern als Architektur-Entscheidungen in konkreten Szenarien.

\subsection{Eval-Pipeline-Design}

\textbf{Frage:} Du hast 500 Golden-Dataset-Fälle. Dein Prompt-Update verbessert 480 Fälle, verschlechtert aber 20. Wie entscheidest du, ob das Update in Produktion geht?

\textbf{System-Review:}

\begin{enumerate}[leftmargin=*]
    \item \textbf{Segment-Analyse}: Die 20 verschlechterten Fälle nach Kategorie clustern (Domain, Intent, Schwierigkeit). Handelt es sich um einen systematischen oder zufälligen Fehler?
    \item \textbf{Produktions-Drift}: Wie viele der 20 Fälle treten in der Produktion tatsächlich auf? Ein Fall, der im Dataset, aber nie in der Realität vorkommt, ist irrelevant.
    \item \textbf{Cost-Impact}: Wie viel kostet der Fehler in der Produktion? Ein Halluzinations-Fehler bei Support-Tickets kostet 5 EUR (Gutschein). Ein Halluzinations-Fehler bei Medikamenten-Dosierung kostet Menschenleben. Der Threshold muss use-case-spezifisch sein.
    \item \textbf{Gradual Rollout}: Das Update in Canary (5\,\% Traffic) ausrollen, Online-Metriken vergleichen, bei positiver Tendenz hochfahren.
\end{enumerate}

\subsection{Produktions-Incident-Response}

\textbf{Frage:} Dein LLM-System zeigt ab 14:23 Uhr plötzlich eine 5-fach höhere Halluzinationsrate. Was ist dein Prozess?

\textbf{Incident-Response:}

\begin{verbatim}
   Schritt 1: Erkennen (Metrik-Alarm)
   Halluzinationsrate > 5% -> PagerDuty-Alarm

   Schritt 2: Eingrenzen (Observability)
   - Trace-Logs checken: Aendern sich die User-Queries um 14:23?
   - Model-Logs checken: Wurde ein Modell-Update deployed?
   - Provider-Status checken: Gibt es ein bekanntes Problem?

   Schritt 3: Isolieren
   - Fallback auf gecachtes Modell (vorherige Version)
   - Traffic auf Backup-Provider umleiten

   Schritt 4: Ursachenanalyse
   - Root Cause: Provider hat "gpt-4o" durch ein neues
     Modell ersetzt (downgrade hinter generischem String)
   - Fix: Datierte Model-Versionen nutzen (gpt-4o-2026-03-15)

   Schritt 5: Praevention
   - Golden-Dataset-Check in CI/CD
   - Provider-Change-Monitoring
   - Modell-Version-Pinning als Standard
\end{verbatim}

% ===========================================================%
% BEST PRACTICES
% ===========================================================%
\section{Best Practices}

\begin{itemize}[leftmargin=*]
    \item \textbf{Zeige Fehler, nicht nur Erfolge}: Ein Repository mit einem Post-Mortem (warum eine Architektur-Entscheidung falsch war) ist überzeugender als eines mit perfektem Code.
    \item \textbf{Messwerte statt Behauptungen}: \glqq Ich habe Kosten gesenkt\grqq{} ist schwach. \glqq Cost per Query von 0,012 EUR auf 0,004 EUR gesenkt durch Semantic Caching\grqq{} ist ein Produktionsnachweis.
    \item \textbf{Jedes Projekt braucht ein README mit Architektur-Diagramm}: Eine ASCII-Skizze der Systemarchitektur mit Datenfluss und Fallback-Pfaden zeigt Systemdenken.
    \item \textbf{CI/CD muss im Repository sichtbar sein}: Ein GitHub-Actions-Workflow, der Tests + Eval + Deployment automatisiert, ist der stärkste Einzelnachweis für Produktionsreife.
    \item \textbf{Bereite zwei Architekturen vor}: Eine für den Idealfall (Standard-Stack) und eine für den Fallback (Degraded Mode). Das zeigt, dass du mit Ausfällen rechnest.
\end{itemize}

% ===========================================================%
% ANTI-PATTERNS
% ===========================================================%
\section{Anti-Patterns}

\begin{itemize}[leftmargin=*]
    \item \textbf{Portfolio ohne Observability}: Ein Repository ohne Logging-Konfiguration, Metriken oder Tracing-Framework beweist keine Produktionserfahrung.
    \item \textbf{Eval-Loop ohne Datensatz}: \glqq Ich teste mit LLM-Judge\grqq{} ohne Golden Dataset ist kein Eval -- es ist eine Meinung. Jeder Eval-Loop braucht versionierte Testfälle.
    \item \textbf{Fallback nur auf dem Papier}: Ein README, das \glqq Fallback: Anthropic\grqq{} erwähnt, aber keinen Code dafür zeigt, ist kein System-Design. Zeige die Fallback-Implementierung.
    \item \textbf{Cache ohne Invalidierung}: Ein Projekt, das Caching erwähnt, aber keine TTL oder Invalidierungsstrategie hat, zeigt mangelndes Systemdenken.
    \item \textbf{Kosten ignorieren}: Projekte, die GPT-4o für alles nutzen (Statusabfragen und komplexe Analysen), zeigen kein Cost-Bewusstsein. Zeige Modell-Routing.
\end{itemize}

% ===========================================================%
% ZUSAMMENFASSUNG
% ===========================================================%
\begin{zusammenfassung}
\begin{itemize}[leftmargin=*]
    \item AI Engineers werden an ihren Systemen gemessen -- nicht an Zertifikaten oder Konzeptwissen. Die sechs Dimensionen: Latenz, Kosten, Zuverlässigkeit, Halluzination, Eval, Observability.
    \item Portfolio-Projekte müssen fünf Layer nachweisen: Kostenkontrolle, Fallback, Eval-Loop, Observability, Deployment.
    \item System-Design-Interviews sind Debugging-Probleme. Behandle jedes Szenario als kaputtes System, das repariert werden muss.
    \item Failure-Mode-Analysen trennen Junior- von Senior-Engineers. Zeige, dass du die fünf häufigsten Produktionsfehler kennst und verhindern kannst.
    \item Ein Repository mit Post-Mortem, Metriken und Fallback-Implementierung ist überzeugender als 100 perfekte Commits ohne Kontext.
\end{itemize}
\end{zusammenfassung}

\begin{merke}
Der stärkste Produktionsnachweis ist nicht das, was du kannst -- es ist das, was du gebaut hast, und wie es unter Druck reagiert. Ein RAG-System mit 99,5\,\% Uptime, 0,003 EUR Cost per Query und dokumentiertem Incident-Response-Plan sagt mehr als jeder Lebenslauf.
\end{merke}

% ===========================================================%
% PRAXISPROJEKT
% ===========================================================%
\section{Praxisprojekt -- Production-Proof-Blueprint}

\textbf{Ziel:} Erstelle ein Repository, das alle fünf Layer eines Production-Proof-Projekts demonstriert.

\textbf{Aufgaben:}
\begin{enumerate}[leftmargin=*]
    \item \textbf{Kostenkontrolle}: Implementiere Semantic Caching + Modell-Routing in einem RAG-System. Dokumentiere die Cost-per-Query-Reduktion.
    \item \textbf{Fallback}: Baue eine Fallback-Chain mit 3 Providern + Degraded Mode. Jeder Ausfall wird geloggt.
    \item \textbf{Eval-Loop}: Erstelle ein Golden Dataset (50 Fälle) und einen CI-Check, der bei $<$ 85\,\% Pass-Rate den Merge blockiert.
    \item \textbf{Observability}: Füge Tracing (OpenTelemetry) hinzu. Erstelle ein Dashboard mit Latenz, Cost, Halluzinationsrate.
    \item \textbf{Deployment}: docker-compose.yml mit 3 Services (API, Redis, Qdrant) + Healthcheck + CI/CD.
\end{enumerate}

\textbf{Erfolgskriterium:} Das Repository kann von einem Senior-Engineer innerhalb von 15 Minuten als produktionsreif eingestuft werden -- nicht weil der Code perfekt ist, sondern weil die Systemstruktur stimmt.

% ============================================================
% INTERVIEWFRAGEN
% ============================================================
\section{Interviewfragen}

\begin{enumerate}[leftmargin=*]
    \item \textbf{System-Design: RAG mit 500 Golden-Dataset Fällen. Update verbessert 480, zerstört 20. Entscheidung?}

    \textbf{Antwort:} Segment-Analyse: Sind die 20 Fälle systematisch oder Zufall? Cost-Impact: Ist ein Fehler dort kritisch (Medizin)? Wenn unkritisch: Canary Rollout (5\%) mit Fallback-Chain.

    \item \textbf{Failure-Mode: RAG hat 12 \% Halluzinationsrate. 5-Schritt Debug?}

    \textbf{Antwort:} 1) Retrieval (Hit Rate). 2) Chunk-Qualität. 3) Prompt-Direktive ("Darf sagen: Weiß ich nicht"). 4) Re-Ranking Effizienz. 5) Golden Dataset Metriken.

    \item \textbf{Multi-Agent kostet plötzlich 4x mehr. Prüfpunkte?}

    \textbf{Antwort:} 1) Step-Count Explusion (Agent hängt im Loop). 2) Token-Verbrauch pro Step (History läuft voll). 3) Modell-Routing kaputt (alles geht an Opus statt Haiku).

    \item \textbf{Accuracy +6 \%, Kosten 5x höher. Entscheidung?}

    \textbf{Antwort:} Reiner Kosten-Nutzen-Case. Ist die +6\% Accuracy im Business-Case 5x mehr Budget wert? Lösung: Dynamisches Routing nur für extrem harte Queries.

    \item \textbf{Provider Rate-Limit, keine Fallback-Infra. Erste 5 Minuten?}

    \textbf{Antwort:} 1) Caches auf max TTL zwingen (Stale Data). 2) Degraded Mode (UI zeigt FAQ). 3) Backup-API-Keys der Entwickler injizieren.

    \item \textbf{"99,9 \% Availability." — Deine Rückfragen an den Provider?}

    \textbf{Antwort:} API-Ping oder E2E Inference? Bezieht sich die SLA auf HTTP 200 oder korrekte, generierte Tokens? Was ist die P99 Time-to-First-Token?

    \item \textbf{Semantic Cache mit 8 \% Hit-Rate. Stellschrauben?}

    \textbf{Antwort:} 1) Confidence-Threshold senken. 2) Text vor Embedding normalisieren (Lowercase/Stopwords). 3) Intent-basiert cachen statt volltextbasiert.

    \item \textbf{Beweise Produktionsreife in 3 Minuten (Artefakte)?}

    \textbf{Antwort:} `docker-compose.yml` (Infra), `.github/workflows/eval.yml` (CI/CD Testing) und ein Screenshot des Datadog/LangSmith Tracing-Dashboards.

    \item \textbf{LLM-System zeigt ab 14:23 5x höhere Halluzinationsrate. Prozess?}

    \textbf{Antwort:} Erkennen (Metrik-Alarm) -> Isolieren (Traffic auf Backup-Provider) -> Ursache finden (Provider hat Modell hinter generischem Alias getauscht) -> Prävention (Version-Pinning).

    \item \textbf{Wie erklärst du einem non-tech Stakeholder, warum das LLM 2+2=5 sagt?}

    \textbf{Antwort:} LLMs rechnen nicht, sie raten das wahrscheinlichste nächste Wort basierend auf Trainingsmustern. Ohne ein Calculator-Tool (Agent) fehlt die mathematische Engine.


\end{enumerate}


% ===========================================================%
% RESSOURCEN
% ===========================================================%
\section{Ressourcen}

\begin{itemize}[leftmargin=*]
    \item \textbf{Incident-Response für ML-Systeme}: \href{https://github.com/opengovernance/ml-incident-response}{github.com/opengovernance/ml-incident-response} -- Runbooks für Produktions-KI-Ausfälle
    \item \textbf{Microsoft AI Debugging}: \href{https://learn.microsoft.com/en-us/azure/machine-learning/how-to-troubleshoot-llm}{learn.microsoft.com} -- Debugging-LLM-Produktionssysteme
    \item \textbf{OpenTelemetry LLM Tracing}: \href{https://opentelemetry.io/docs/specs/otel/overview/}{opentelemetry.io} -- Standard für LLM-Observability
    \item \textbf{Production ML Monitoring}: \href{https://www.evidentlyai.com/}{evidentlyai.com} -- ML-Monitoring und Drift-Erkennung
    \item \textbf{System Design Interview (Alex Xu)}: Band 1 + 2 -- übertragbar auf AI-System-Design: Chat-Systeme, Data-Pipelines, Rate-Limiter
    \item \textbf{Failure Mode Analysis}: \href{https://github.com/mitre/atlas}{MITRE ATLAS} -- Taktiken für AI-System-Angriffe und Ausfälle
\end{itemize}

\section*{Abschluss}

AI Engineering ist System-Engineering. Der Unterschied zwischen einem guten Notebook und einem guten AI Engineer ist die Fähigkeit, unter Druck zu liefern: wenn die Halluzinationsrate steigt, wenn der Provider ausfällt, wenn die Kosten explodieren. Die Systeme, die du baust, müssen das aushalten -- und deine Projekte müssen beweisen, dass sie es können.

Der nächste Schritt ist nicht \glqq bewirb dich\grqq{}. Der nächste Schritt ist: \glqq Baue ein System, das unter Druck nicht versagt. Und dann zeig es.\grqq
